\chapter*{\rm \large \bf ABSTRACT}
\vspace{3.0mm}
\begin{justify}

\hspace{0.5cm}Rust has emerged as a prominent systems programming language due to its strong guarantees of memory safety, concurrency, and high-performance execution. Despite these compile-time advantages, the operational quality of Rust applications remains highly dependent on localized developer implementation practices, including fine-grained ownership boundaries, micro-allocation behaviors, and loop safety mechanics. Existing static analysis ecosystems operate primarily on linear token patterns, providing fragmented, rule-based diagnostic alerts devoid of contextual reasoning regarding localized lifecycle mutations. Furthermore, current infrastructure tools lack the cross-layer evaluation models necessary to reason about how a localized memory allocation anti-pattern simultaneously degrades performance and code maintainability indicators inside a single subroutine body.
\par This paper presents RustAuditAI, an intra-procedural software quality intelligence framework that performs multi-dimensional semantic graph analysis bounded within individual Rust functions. The proposed system parses localized function syntax and constructs deep, interconnected semantic abstractions, including functional ASTs, localized CFGs, and dedicated FLOGs. By executing specialized traversal algorithms across this localized Code Property Graph structure, RustAuditAI unifies structural metrics with semantic analysis across four primary quality vectors: memory safety verification, code execution complexity, localized resource density, and defensive security. These analysis vectors are dynamically synthesized via a weighted valuation model to compute a standardized Rust Quality Index (RQI) representing localized code health. Bounded by the deterministic structural properties of the generated semantic graphs, an explainable recommendation engine generates context-aware, human-interpretable refactoring suggestions that specify the underlying structural degradation paths and estimate the quantitative execution optimizations achieved through more idiomatic Rust idioms.
\end{justify}